\begin{itemize}
\item Die Wärme entsteht bei der Verbrennung im Zylinder. Durch \emph{Wärmeleitung} gelangt sie durch die Zylinderwände in die Kanäle der Kühlflüssigkeit.
\item Die Wasserpumpe treibt die Kühlflüssigkeit zum Kühler (Radiator): Die Wärme wird mit der Flüssigkeit mitgeführt, also durch \emph{erzwungene Konvektion} transportiert.
\item Im Kühler gelangt die Wärme durch \emph{Wärmeleitung} durch die dünnen Metallwände und Lamellen an die Oberfläche. Die grosse Oberfläche gibt sie an die vorbeiströmende Luft ab, die vom Fahrtwind oder vom Ventilator weggeführt wird (\emph{Konvektion}).
\item Ein grosser Teil der Abwärme verlässt den Motor direkt mit den heissen Abgasen durch den Auspuff (Konvektion).
\item Motorblock, Auspuff und Kühler geben zudem einen kleineren Teil durch \emph{Wärmestrahlung} und durch Konvektion der umgebenden Luft ab.
\end{itemize}
